\chapter{Classical Linear Model}

\section{Model and Assumptions}

\[
y_t = x_t' \beta + e_t, \quad t=1,\dots,T
\]

Matrix form:
\[
y = X\beta + e
\]

\subsection{Assumptions}

\begin{itemize}
    \item[A1] Linearity:
    \[
    E(y_t) = x_t' \beta
    \]
    \item[A2] $X$ is non-stochastic
    \item[A3] $e_t$ independent
    \item[A4] Homoskedasticity:
    \[
    V(e_t)=\sigma^2
    \]
    \item[A5] Full rank:
    \[
    \det(X'X) \neq 0
    \]
\end{itemize}

\section{Least Squares Estimation}

\[
S = e'e = (y - X\beta)'(y - X\beta)
\]

First order condition:
\[
\frac{\partial S}{\partial \beta} = -2X'y + 2X'X\beta = 0
\]

OLS estimator:
\[
\hat{\beta} = (X'X)^{-1}X'y
\]

\section{Maximum Likelihood Estimation}

Assume:
\[
e_t \sim N(0,\sigma^2)
\]

Log-likelihood:
\[
\ell = -\frac{T}{2}\ln(2\pi) - \frac{T}{2}\ln\sigma^2 - \frac{(y-X\beta)'(y-X\beta)}{2\sigma^2}
\]

MLE:
\[
\hat{\beta}_{ML} = (X'X)^{-1}X'y
\]

\[
\hat{\sigma}^2 = \frac{(y-X\hat{\beta})'(y-X\hat{\beta})}{T}
\]

\section{Properties of $\hat{\beta}$}

\subsection{Unbiasedness}

\[
E(\hat{\beta}) = \beta
\]

\subsection{Variance}

\[
V(\hat{\beta}) = \sigma^2 (X'X)^{-1}
\]

\subsection{Gauss-Markov Theorem}

$\hat{\beta}$ is BLUE.

\section{Variance Estimator}

\[
\hat{\sigma}^2 = \frac{(y-X\hat{\beta})'(y-X\hat{\beta})}{T}
\]

Biased:
\[
E(\hat{\sigma}^2) = \sigma^2 \frac{T-K}{T}
\]

Unbiased:
\[
s^2 = \frac{(y-X\hat{\beta})'(y-X\hat{\beta})}{T-K}
\]

\section{Large Sample Properties}

\subsection{Consistency}

\[
\text{plim } \hat{\beta} = \beta
\]

\subsection{Asymptotic Variance}

\[
V(\hat{\beta}) \approx \frac{\sigma^2}{T}Q^{-1}
\]

\subsection{Asymptotic Normality}

\[
\hat{\beta} \sim N\left(\beta, s^2 (X'X)^{-1}\right)
\]

\section{Prediction}

\[
y_0 = X_0 \beta + e_0
\]

Prediction:
\[
\hat{y}_0 = X_0 \hat{\beta}
\]

\subsection{Prediction Variance}

\[
V(y_0 - \hat{y}_0) = \sigma^2 \left[1 + X_0 (X'X)^{-1} X_0'\right]
\]

\section{Goodness of Fit}

\[
SST = \sum (y_t - \bar{y})^2
\]

\[
SSE = \sum (y_t - x_t'\hat{\beta})^2
\]

\[
SSR = SST - SSE
\]

\subsection{$R^2$}

\[
R^2 = 1 - \frac{SSE}{SST}
\]

\subsection{Adjusted $R^2$}

\[
\bar{R}^2 = 1 - \frac{SSE/(T-K)}{SST/(T-1)}
\]



\section{Restricted Least Squares}

Consider the linear restrictions
\[
R\beta = r,
\]
where $R$ is a $J \times K$ matrix with rank $J$, and $r$ is a $J \times 1$ vector.

The restricted least squares estimator solves
\[
\min_{\beta} (y-X\beta)'(y-X\beta)
\quad \text{subject to} \quad R\beta = r.
\]

The Lagrangian is
\[
\mathcal{L}
=
(y-X\beta)'(y-X\beta)
+
\lambda'(r-R\beta).
\]

The first order condition is
\[
-2X'y + 2X'X\beta - R'\lambda = 0.
\]

Thus,
\[
\beta
=
(X'X)^{-1}X'y
+
(X'X)^{-1}R'\frac{\lambda}{2}.
\]

Since
\[
\hat{\beta} = (X'X)^{-1}X'y,
\]
we have
\[
\beta
=
\hat{\beta}
+
(X'X)^{-1}R'\frac{\lambda}{2}.
\]

Using the restriction $R\beta = r$,
\[
r
=
R\hat{\beta}
+
R(X'X)^{-1}R'\frac{\lambda}{2}.
\]

Therefore,
\[
\frac{\lambda}{2}
=
\left[R(X'X)^{-1}R'\right]^{-1}
(r-R\hat{\beta}).
\]

Hence, the restricted estimator is
\[
\hat{\beta}_R
=
\hat{\beta}
+
(X'X)^{-1}R'
\left[R(X'X)^{-1}R'\right]^{-1}
(r-R\hat{\beta}).
\]

\section{Properties of the Restricted Estimator}

\subsection{Unbiasedness}

Taking expectations,
\[
E(\hat{\beta}_R)
=
\beta
+
(X'X)^{-1}R'
\left[R(X'X)^{-1}R'\right]^{-1}
(r-R\beta).
\]

Therefore,
\[
E(\hat{\beta}_R)=\beta
\quad \text{if} \quad R\beta=r.
\]

That is, $\hat{\beta}_R$ is unbiased if the imposed restrictions are true.

If the restrictions are false, then
\[
E(\hat{\beta}_R) \neq \beta.
\]

\subsection{Variance}

Define
\[
M^*
=
I_K
-
(X'X)^{-1}R'
\left[R(X'X)^{-1}R'\right]^{-1}R.
\]

Then,
\[
\hat{\beta}_R
=
M^*\hat{\beta}
+
(X'X)^{-1}R'
\left[R(X'X)^{-1}R'\right]^{-1}r.
\]

Since the second term is non-random,
\[
V(\hat{\beta}_R)
=
M^* V(\hat{\beta}) M^{*'}.
\]

Using
\[
V(\hat{\beta}) = \sigma^2 (X'X)^{-1},
\]
we obtain
\[
V(\hat{\beta}_R)
=
\sigma^2
\left[
(X'X)^{-1}
-
(X'X)^{-1}R'
\left[R(X'X)^{-1}R'\right]^{-1}
R(X'X)^{-1}
\right].
\]

Therefore,
\[
V(\hat{\beta}_R) \leq V(\hat{\beta}).
\]

The restricted estimator has smaller variance than the unrestricted estimator, but this advantage is meaningful only when the restrictions are true.

\section{Hypothesis Testing with Linear Restrictions}

Consider the hypothesis
\[
H_0: R\beta = r,
\]
against
\[
H_1: R\beta \neq r.
\]

Here, $J$ is the number of restrictions.

Let
\[
SSE_u = (y-X\hat{\beta})'(y-X\hat{\beta})
\]
be the unrestricted sum of squared errors, and
\[
SSE_R = (y-X\hat{\beta}_R)'(y-X\hat{\beta}_R)
\]
be the restricted sum of squared errors.

Since the restricted model imposes additional constraints,
\[
SSE_R \geq SSE_u.
\]

\subsection{F-Test}

The test statistic is
\[
F
=
\frac{(SSE_R-SSE_u)/J}{SSE_u/(T-K)}.
\]

Under the null hypothesis,
\[
F \sim F(J,T-K).
\]

The testing rule is:
\[
\text{Reject } H_0
\quad \text{if} \quad
F > F_{\alpha}(J,T-K).
\]

Equivalently, using the unrestricted estimator,
\[
F
=
\frac{
(R\hat{\beta}-r)'
\left[R(X'X)^{-1}R'\right]^{-1}
(R\hat{\beta}-r)/J
}{
s^2
},
\]
where
\[
s^2 = \frac{SSE_u}{T-K}.
\]

\section{Special Case: One Linear Restriction}

Suppose there is only one restriction, so that $J=1$.

The null hypothesis is
\[
H_0: R\beta = r.
\]

The test statistic becomes
\[
t
=
\frac{R\hat{\beta}-r}
{
s\sqrt{R(X'X)^{-1}R'}
}.
\]

Under $H_0$,
\[
t \sim t(T-K).
\]

Because
\[
t^2 \sim F(1,T-K),
\]
the one-restriction $t$-test is equivalent to the corresponding $F$-test.

\section{Testing a Single Coefficient}

A common special case is
\[
H_0: \beta_i = r.
\]

Then the test statistic is
\[
t
=
\frac{\hat{\beta}_i-r}{se(\hat{\beta}_i)}.
\]

Under the null hypothesis,
\[
t \sim t(T-K).
\]

The confidence interval for $\beta_i$ is
\[
\hat{\beta}_i
\pm
t_{\alpha/2,T-K} \, se(\hat{\beta}_i).
\]

If $r$ lies outside this interval, then $H_0$ is rejected.

\section{Likelihood Ratio Test}

Under normality,
\[
e \sim N(0,\sigma^2 I_T).
\]

The likelihood ratio statistic compares the unrestricted and restricted maximized likelihoods:
\[
\lambda
=
\frac{L_u^*}{L_R^*}.
\]

Since the restricted model cannot fit better than the unrestricted model,
\[
\lambda \geq 1.
\]

Using the maximum likelihood variance estimators,
\[
\hat{\sigma}_u^2
=
\frac{SSE_u}{T},
\qquad
\hat{\sigma}_R^2
=
\frac{SSE_R}{T},
\]
we obtain
\[
\lambda
=
\left(
\frac{\hat{\sigma}_R^2}{\hat{\sigma}_u^2}
\right)^{T/2}.
\]

Since
\[
\frac{\hat{\sigma}_R^2}{\hat{\sigma}_u^2}
=
\frac{SSE_R}{SSE_u},
\]
the likelihood ratio test is closely related to the $F$-test.

\section{Asymptotic Tests}

For large samples, we can use asymptotic distributions.

Under $H_0$,
\[
\frac{
(R\hat{\beta}-r)'
\left[R(X'X)^{-1}R'\right]^{-1}
(R\hat{\beta}-r)
}{
s^2
}
\sim
\chi^2(J).
\]

Thus, the asymptotic Wald test statistic is
\[
W
=
\frac{
(R\hat{\beta}-r)'
\left[R(X'X)^{-1}R'\right]^{-1}
(R\hat{\beta}-r)
}{
s^2
}.
\]

The decision rule is
\[
\text{Reject } H_0
\quad \text{if} \quad
W > \chi^2_{\alpha}(J).
\]

For a single restriction,
\[
\frac{R\hat{\beta}-r}
{
s\sqrt{R(X'X)^{-1}R'}
}
\approx N(0,1).
\]

Therefore, in large samples, the $t$-test can be approximated by a standard normal test.